\chapter{Un projet de A à Z, en cinq commandes}
\label{ch:aaz}

Ce chapitre part d'un dossier vide et arrive à un programme qui s'exécute. Cinq
commandes, et rien d'autre.

\textbf{Toutes les sorties de ce chapitre sont des captures réelles}, prises en
enchaînant ces commandes. Rien n'est recopié de mémoire.

\section{Commande 1 --- créer l'espace de travail}

\begin{nkcode}{}
jenga workspace Bonjour --no-interactive
\end{nkcode}

Trois fichiers apparaissent :

\begin{nkcode}{Ce que la commande a produit}
Bonjour/.gitignore
Bonjour/Bonjour.jenga
Bonjour/pyrightconfig.json
\end{nkcode}

\begin{nkretenir}
\textbf{Un des trois ne sert qu'à votre éditeur.}
\texttt{pyrightconfig.json} fait que l'auto-complétion connaît le vocabulaire
Jenga dans vos \texttt{.jenga} --- qui sont du Python, donc analysables. Il
n'appartient à aucun éditeur en particulier : tout client \textsc{lsp} le lit,
NKCode compris.

Le seul fichier qui compte pour la construction est \texttt{Bonjour.jenga}.
\end{nkretenir}

\begin{nkpiege}
\textbf{Jusqu'au 2 septembre 2026, il y en avait \emph{quatre}} --- un
\texttt{.vscode/settings.json} apparaissait aussi, pour un éditeur que vous
n'aviez peut-être jamais ouvert. La fonction de détection finissait par
« si rien n'est détecté, on assume VSCode ».

Ce n'est pas le dossier de trop qui était grave, c'est qu'il \textbf{se
confirmait lui-même} : le \texttt{.vscode/} ainsi créé était \emph{ensuite}
détecté par le test qui décide d'écrire, donc une \emph{supposition} du premier
passage devenait un \emph{fait} au second, définitivement --- et plus personne
ne pouvait répondre à « qui a voulu ce dossier ? ».

Jenga n'écrit désormais que pour un éditeur \textbf{dont la trace existe} :
\texttt{.vscode/}, \texttt{.idea/}, \texttt{.nkcode/}\ldots Sans trace, seul
\texttt{pyrightconfig.json} sort. Un utilisateur de VSCode qui part de zéro le
demande une fois, et c'est tout :

\begin{nkcode}{}
jenga ide-setup --editor vscode
\end{nkcode}

\emph{Un opt-in explicite vaut mieux qu'un dossier apparu tout seul.}
\end{nkpiege}

Et voici ce qu'il contient --- lisez-le, il tient en dix lignes utiles :

\begin{nkcode}{Bonjour/Bonjour.jenga, genere}
from Jenga import *

with workspace("Bonjour"):
    configurations(['Debug', 'Release'])
    targetoses([TargetOS.WINDOWS, TargetOS.LINUX, TargetOS.MACOS])
    targetarchs([TargetArch.X86_64])

    # Default toolchain (auto-detected)
    # usetoolchain("host-gcc")

    # Add your projects here
    # with project("MyApp"):
    #     consoleapp()
    #     language("C++")
    #     files(["src/**.cpp"])
\end{nkcode}

\begin{nkretenir}
\textbf{Aucun projet encore, et c'est délibéré.} Un workspace décrit le
\emph{cadre} : quelles configurations existent, quels systèmes on vise, quelles
architectures. Les projets viennent ensuite --- et le gabarit vous montre déjà
leur forme, en commentaire.
\end{nkretenir}

\section{Commande 2 --- créer le projet}

\begin{nkcode}{}
jenga project Bonjour --kind console --no-interactive
\end{nkcode}

\begin{nkpiege}[Le nom du \emph{kind} n'est pas celui de la fonction]
En ligne de commande, c'est \texttt{console}. Dans le \texttt{.jenga}, la
fonction s'appelle \texttt{consoleapp()}. Se tromper donne, textuellement :

\begin{nkcode}{Capture reelle, code de sortie 2}
jenga project: error: argument --kind: invalid choice: 'consoleapp'
(choose from console, windowed, static, shared, test)
\end{nkcode}

\textbf{Le message est bon} : il énumère les cinq valeurs admises. C'est le
genre d'erreur qui se résout en la lisant --- et qui fait perdre dix minutes
quand on ne la lit pas.
\end{nkpiege}

La commande crée la source et \textbf{modifie le workspace} :

\begin{nkcode}{Bonjour/Bonjour.jenga, apres `jenga project`}
    # Project: Bonjour
    with project("Bonjour"):
        consoleapp()
        language("C++")
        location("Bonjour")
        files(["src/**.cpp", "include/**.hpp"])
\end{nkcode}

\begin{nkcode}{Bonjour/Bonjour/src/main.cpp, genere}
#include <iostream>

int main() {
    std::cout << "Hello from Bonjour!" << std::endl;
    return 0;
}
\end{nkcode}

\begin{nkretenir}
\textbf{Cinq lignes décrivent entièrement un programme.} Son genre, son langage,
où il vit, et quels fichiers lui appartiennent. Tout le reste --- compilateur,
drapeaux, dossier de sortie --- a une valeur par défaut que vous ne redéfinirez
que si elle ne vous convient pas.
\end{nkretenir}

\section{Commande 3 --- regarder ce que Jenga a compris}

\begin{nkcode}{jenga info -- capture reelle, abregee}
=========================== Jenga Workspace: Bonjour ===========================

Configurations: Debug, Release
Platforms: Windows
Target OSes: Windows, Linux, macOS
Target Architectures: x86_64

Projects
------------------------------------------------------------
Name      Kind         Language   Test   External
=================================================
Bonjour   ConsoleApp   C++        No     No

Available Toolchains
------------------------------------------------------------
Name                 Family        Target OS   Arch     Env
===============================================================
host-clang           clang         Windows     x86_64   mingw
host-gcc             gcc           Windows     x86_64   mingw
msvc                 msvc          Windows     x86_64   msvc
clang-mingw          clang         Windows     x86_64   mingw
mingw                gcc           Windows     x86_64   mingw
clang-cross-linux    clang         Linux       x86_64   gnu
android-ndk          android-ndk   Android     arm64    android
emscripten           emscripten    Web         wasm32
\end{nkcode}

\begin{nkretenir}
\textbf{\texttt{jenga info} est la première commande à taper quand quelque chose
surprend.} Elle montre ce que Jenga a \emph{compris} de votre description ---
pas ce que vous croyez avoir écrit.

Notez la distinction entre \emph{Platforms: Windows} (la machine sur laquelle on
est) et \emph{Target OSes: Windows, Linux, macOS} (ce que le workspace déclare
viser). Confondre les deux fait chercher pourquoi « Linux n'est pas là » alors
qu'il l'est.

Et la table des chaînes d'outils est \textbf{ce qui est disponible ici, sur cette
machine} : Android et Emscripten y figurent parce que leurs SDK sont installés.
Sur une autre machine, la table est plus courte.
\end{nkretenir}

\section{Commande 4 --- construire}

\begin{nkcode}{jenga build --config Release -- capture reelle}
Loading workspace...

Configuration: Release
Target:        Windows x86_64
Toolchain:     clang-mingw

Build Order (1 projects):
  1. Bonjour [CONSOLE_APP]

  Project: Bonjour                                     Kind: CONSOLE_APP

Found 1 source file(s)
  [1/1] Compiled: main.cpp
Linking...
Built: Build\Bin\Release-Windows\Bonjour\Bonjour.exe

  Build Successful                                        Time: 0.83s

                                BUILD COMPLETED
Projects Built:  1/1
Time:           0.83s
Status:         SUCCESS
\end{nkcode}

\begin{nkretenir}
\textbf{Quatre lignes de cette sortie méritent d'être lues à chaque fois}, et ce
sont celles qu'on saute :

\begin{center}
\begin{tabular}{@{}ll@{}}
\toprule
\textbf{Ligne} & \textbf{Ce qu'elle évite de croire à tort} \\
\midrule
\texttt{Configuration: Release} & que vous mesurez du Debug \\
\texttt{Toolchain: clang-mingw} & que vous compilez avec MSVC \\
\texttt{Build Order} & qu'un projet est construit alors qu'il ne l'est pas \\
\texttt{Compiled: main.cpp} & qu'un fichier a été recompilé \\
\bottomrule
\end{tabular}
\end{center}

La dernière est la plus importante, et le chapitre 10 lui consacre une
démonstration : \textbf{un build vert prouve que ce qu'il a construit compile,
pas que ce que vous avez écrit a été construit.}
\end{nkretenir}

\begin{nkretenir}[Le chemin de sortie se lit, il ne se devine pas]
\texttt{Build\textbackslash Bin\textbackslash Release-Windows\textbackslash Bonjour\textbackslash Bonjour.exe}

La forme est \texttt{Build/Bin/<config>-<système>/<projet>/}. Elle change avec la
configuration \emph{et} avec le système --- et sous Linux elle porte en plus le
backend graphique, par exemple \texttt{Release-Linux-xlib}. Un script qui code ce
chemin en dur échoue sur une construction \emph{réussie}.
\end{nkretenir}

\section{Commande 5 --- lancer}

\begin{nkcode}{jenga run --config Release -- capture reelle}
  EXECUTION  --  Bonjour.exe
     ...\Build\Bin\Release-Windows\Bonjour\Bonjour.exe

Application console sans terminal : ouverture d'une console dediee.
\end{nkcode}

\begin{nkretenir}
\texttt{jenga run} retrouve l'exécutable tout seul, à partir de la configuration
et de la plateforme. Vous n'avez pas à connaître le chemin --- ce qui est
justement pourquoi il faut le lire au moins une fois, à l'étape précédente.

\textbf{Et il ouvre une console dédiée} quand l'application est de type console
et qu'aucun terminal n'est attaché. C'est une commodité ; elle explique aussi
pourquoi la sortie du programme n'apparaît pas toujours dans \emph{votre}
terminal.
\end{nkretenir}

\section{Le récapitulatif, en cinq lignes}

\begin{nkcode}{De rien a un programme qui tourne}
jenga workspace Bonjour --no-interactive
cd Bonjour
jenga project Bonjour --kind console --no-interactive
jenga build --config Release
jenga run   --config Release
\end{nkcode}

\begin{nkretenir}
\textbf{Aucune de ces commandes n'a demandé de compilateur, de chemin, ni de
drapeau.} C'est le contrat de Jenga : vous décrivez ce que vous voulez, il
choisit comment.

Le reste du livre porte sur les moments où ce choix ne vous convient pas --- et
sur la façon d'en reprendre la main sans perdre le reste.
\end{nkretenir}

\begin{nkbilan}
Cinq commandes suffisent : \texttt{workspace}, \texttt{project}, \texttt{info},
\texttt{build}, \texttt{run}. Le workspace décrit le cadre, le projet décrit un
artefact en cinq lignes, et \texttt{info} montre ce que Jenga a \emph{compris} --
la première commande à taper quand quelque chose surprend. La sortie de
construction porte quatre lignes qu'il faut lire à chaque fois : la
configuration, la chaîne d'outils, l'ordre, et les fichiers réellement
recompilés. Enfin, le chemin de sortie dépend de la configuration \emph{et} du
système : il se lit, il ne se code pas en dur.
\end{nkbilan}

\begin{nkquestions}
\begin{enumerate}
\item Quels fichiers \texttt{jenga workspace} produit-il, et lequel compte
      vraiment ?
\item Pourquoi \texttt{-{}-kind consoleapp} échoue-t-il alors que
      \texttt{consoleapp()} existe ?
\item Quelle différence entre \emph{Platforms} et \emph{Target OSes} dans la
      sortie de \texttt{info} ?
\item Citez les quatre lignes de la sortie de \texttt{build} qu'il faut lire, et
      ce que chacune évite.
\item De quoi dépend le chemin de l'exécutable produit ?
\end{enumerate}
\end{nkquestions}

\begin{nkpratique}
Refaites les cinq commandes, dans un dossier vide, sans regarder ce chapitre.

Puis trois variantes, chacune répondant à une question :
\begin{enumerate}
\item construisez en \texttt{Debug} --- où l'exécutable atterrit-il ?
\item changez le message de \texttt{main.cpp} et reconstruisez --- combien de
      fichiers sont recompilés ?
\item lancez \texttt{jenga build} \emph{sans} modifier quoi que ce soit ---
      qu'affiche la ligne \texttt{Compiled:} ?
\end{enumerate}
La troisième prépare le chapitre 10.
\end{nkpratique}

\begin{nkdemo}
Prouvez que \texttt{jenga info} lit votre description, et non un état mis en
cache.

Ajoutez un second projet à votre workspace --- cinq lignes suffisent --- et
relancez \texttt{jenga info} \textbf{sans construire}. Le tableau des projets
doit en montrer deux.

Retirez-le, relancez : un seul.

\textbf{Ce que la démonstration établit}, et c'est ce qui la rend utile : la
description est relue \emph{à chaque commande}. Vous pouvez donc vous fier à
\texttt{info} pour vérifier une modification avant de payer une construction.
\end{nkdemo}

\begin{nklogique}
Vous lancez \texttt{jenga build}, la sortie dit \texttt{SUCCESS}, et pourtant
votre programme se comporte comme avant.

\begin{enumerate}
\item Énumérez au moins quatre explications, en distinguant celles où la
      compilation \emph{n'a pas eu lieu} de celles où elle a eu lieu \emph{sans
      effet}.
\item Pour chacune, dites quelle ligne de la sortie de \texttt{build} la
      confirmerait ou l'éliminerait.
\item Une de ces explications ne se voit dans \textbf{aucune} ligne de cette
      sortie. Laquelle, et par quoi la prendriez-vous ?
\end{enumerate}
\end{nklogique}
